\documentclass[11pt,a4paper]{article}
\usepackage[T1]{fontenc}
\usepackage[utf8]{inputenc}
\usepackage[main=italian,provide=*]{babel}
\usepackage{amsmath,amssymb}
\usepackage{geometry}
\geometry{margin=2.3cm,top=2.5cm,bottom=2.7cm}
\usepackage{booktabs}
\usepackage[table]{xcolor}
\usepackage{tcolorbox}
\tcbuselibrary{skins,breakable}
\usepackage{titlesec}
\usepackage{enumitem}
\usepackage{needspace}
\usepackage{fancyhdr}
\usepackage{hyperref}
\hypersetup{colorlinks=true,linkcolor=Sepia,citecolor=Sepia,urlcolor=Sepia}
\newcommand{\Tr}{\operatorname{Tr}}
\newcommand{\ket}[1]{\lvert #1 \rangle}
\newcommand{\bra}[1]{\langle #1 \rvert}

% --- palette ---
\definecolor{inkblue}{HTML}{1B3A5C}
\definecolor{lightblue}{HTML}{EAF1F8}
\definecolor{accent}{HTML}{B0562A}
\definecolor{softgray}{HTML}{6B6B6B}
\definecolor{okgreen}{HTML}{2E6B3E}
\definecolor{okgreenbg}{HTML}{EAF5EC}

% --- titoli di sezione ---
\titleformat{\section}
  {\normalfont\Large\bfseries\color{inkblue}}
  {\thesection}{0.6em}{}[{\color{inkblue!40}\titlerule}]
\titlespacing{\section}{0pt}{0.7em}{0.4em}
\titleformat{\subsection}
  {\normalfont\large\bfseries\color{accent}}
  {}{0em}{}
\titlespacing{\subsection}{0pt}{0.5em}{0.2em}

\pagestyle{fancy}
\fancyhf{}
\renewcommand{\headrulewidth}{0.4pt}
\fancyhead[L]{\small\color{softgray}Trotter sotto rumore --- Passo 3, Parte 2}
\fancyhead[R]{\small\color{softgray}\thepage}
\fancyfoot[C]{\small\color{softgray}Tesi triennale, Samuele --- Relatore: Prof.\ Alessandro Chiesa}

% --- box riutilizzabili ---
\newtcolorbox{keyeq}[1][]{
  colback=lightblue, colframe=inkblue!70, boxrule=0.6pt,
  arc=2pt, left=8pt, right=8pt, top=3pt, bottom=3pt, #1
}
\newtcolorbox{okbox}[1]{
  colback=okgreenbg, colframe=okgreen!60, boxrule=0.6pt,
  arc=2pt, left=7pt, right=7pt, top=2pt, bottom=2pt,
  fonttitle=\bfseries\color{okgreen}, title={#1}
}

\title{\vspace{-1.5em}\color{inkblue}Quantum simulation (Trotter) sotto rumore --- Passo 3\\[0.2em]
\Large\normalfont\color{softgray}Prima comparsa concreta del compromesso $a/N+g\varepsilon_{2q}N$}
\author{Note di lavoro --- Tesi triennale, Samuele\\ Relatore: Prof.\ Alessandro Chiesa}
\date{}

\begin{document}
\sloppy
\maketitle
\thispagestyle{fancy}
\vspace{-2.2em}

File: \texttt{trotter\_rumoroso\_dimero.py}, \texttt{validate\_trotter\_rumoroso\_dimero.py}.
Punto di lavoro: $b/J=0.35$, $D/J=0.80$, $J=1$, $t=2$.

\section*{Metodologia}
\addcontentsline{toc}{section}{Metodologia}

\begin{itemize}[leftmargin=1.3em,itemsep=0.3em]
\item \textbf{Transpilazione per blocco} (decisione della sessione
precedente, qui applicata per la prima volta): il singolo passo di Trotter
è transpilato \emph{una sola volta} al suo costo minimo (3 CNOT), poi il
blocco già compilato è composto $N$ volte con \texttt{QuantumCircuit.compose}
--- mai ritranspilato per intero. Verificato esplicitamente: ritranspilare
il circuito completo a livello alto ricade nella stessa trappola già
trovata (collasso a costo costante, indipendente da $N$).
\item \textbf{Osservabile}: fedeltà rispetto allo stato bersaglio
\emph{fisico} --- l'evoluzione esatta e continua (nessun Trotter, nessun
rumore) del ground state esatto,
\begin{keyeq}
\centering
$\ket{\psi_\text{target}(t)}=e^{-iHt}\ket{\psi_0^\text{esatto}}, \qquad
F(N)=\bra{\psi_\text{target}(t)}\rho_N\ket{\psi_\text{target}(t)}$
\end{keyeq}
dove $\rho_N$ è lo stato dopo preparazione VQE (Passo 2) + $N$ passi di
Trotter rumorosi. Un'unica quantità impacchetta tutti e tre i contributi
già caratterizzati nella sessione precedente: preparazione (costante in
$N$), Trotter (decresce con $N$), rumore accumulato (cresce con $N$).
\end{itemize}

\section*{Verifica}
\addcontentsline{toc}{section}{Verifica}

\begin{okbox}{Tre controlli superati}
\begin{itemize}[leftmargin=1.2em,itemsep=0.2em]
\item \textbf{Conteggio gate}: $\text{CNOT totali}=2+3N$ (2 dall'ansatz,
Passo 2) confermato esattamente per $N=1,\dots,40$ --- crescita lineare,
nessuna fusione dei passi.
\item \textbf{Cross-check indipendente} (rumore nullo): fedeltà dal
circuito Qiskit vs. potenza di matrice in \texttt{numpy} puro (nessun
circuito) --- coincidono entro $4\times10^{-7}$, compatibile con
l'infedeltà di preparazione già nota (Passo 2), non rumore numerico.
\item \textbf{$N^*$ a un valore finito}, non a $N\to\infty$: prima
comparsa concreta del compromesso predetto teoricamente.
\end{itemize}
\end{okbox}

\needspace{18\baselineskip}
\section*{Risultato: fedeltà in funzione di $N$}
\addcontentsline{toc}{section}{Risultato}

\begin{center}
\renewcommand{\arraystretch}{1.2}
\begin{tabular}{@{}ccc@{}}
\toprule
$N$ & $F$ (rumore nullo) & $F$ (rumore \texttt{ibm\_torino}) \\
\midrule
$1$ & $0.7799372$ & $0.7567291$ \\
$2$ & $0.4776280$ & $0.4624084$ \\
$5$ & $0.8323751$ & $0.7560090$ \\
$8$ & --- & $0.8177253$ \quad ($\approx N^*$) \\
$10$ & $0.9666876$ & $0.8087205$ \\
$20$ & $0.9921835$ & $0.7097565$ \\
$40$ & $0.9980771$ & $0.5427382$ \\
$80$ & $0.9995212$ & $0.3567462$ \\
$160$ & $0.9998804$ & $0.2686760$ \\
\bottomrule
\end{tabular}
\end{center}

\noindent Senza rumore, $F(N)\to1$ (limitato solo dall'infedeltà di
preparazione, trascurabile) al crescere di $N$: puro errore di Trotter, in
diminuzione monotona per $N\gtrsim5$ (la non-monotonia fra $N=1$ e $N=2$ a
$t=2$ non è un errore --- a $N$ piccolo $\tau=t/N$ non è nel regime
perturbativo dove vale $O(1/N^2)$).

\begin{okbox}{$N^*$ individuato}
Con il rumore di riferimento, $F(N)$ ha un \textbf{massimo a $N^*=8$},
$F(N^*)=0.8177$ --- non a $N\to\infty$. Per $N<N^*$ domina l'errore di
Trotter (non ancora abbastanza passi); per $N>N^*$ domina il rumore
accumulato ($2+3N$ CNOT, ciascuno con probabilità $\lambda_{2q}$ di errore).
\end{okbox}

\needspace{6\baselineskip}
\section*{Prossimo passo}
\addcontentsline{toc}{section}{Prossimo passo}

Passo 4: correlazioni dinamiche sotto rumore (circuito Hadamard test
completo, ancilla inclusa, con lo stesso blocco Trotter già compilato).
Poi Passo 5: lo scan vero e proprio su
$(\varepsilon_{1q},\varepsilon_{2q},p_\text{readout})$ per vedere come si
sposta $N^*$ --- qui misurato a un solo punto di riferimento.

\end{document}
